\section{Toolkit and Diagnostics}

\tool treats a tokenizer export as an artifact, not as an opaque component of a
trained recommender. The minimum input is a table
\texttt{sid\_assignments(item\_id, sid\_0, ..., sid\_L)}. Optional tables
provide item metadata, interaction histories, and generator outputs. This
contract lets adapters normalize heterogeneous repositories while keeping the
audit independent of a particular training loop.

\begin{figure}[t]
\centering
\scriptsize
\setlength{\fboxsep}{3pt}
\fbox{%
\begin{minipage}{0.94\columnwidth}
\textbf{Tokenizer exports}:
repo SIDs, public mappings, or controlled sanity tokenizers

\vspace{2pt}\centering$\downarrow$\vspace{2pt}

\raggedright
\textbf{Artifact contract}:
\texttt{sid\_assignments}, \texttt{item\_metadata}, \texttt{interactions},
optional \texttt{generator\_outputs}

\vspace{2pt}\centering$\downarrow$\vspace{2pt}

\raggedright
\textbf{Validation}:
join coverage, prefix depth, provenance, and adapter caveats

\vspace{2pt}\centering$\downarrow$\vspace{2pt}

\raggedright
\textbf{Diagnostics}:
D1 utilization; D2 collisions; D3 collaborative alignment; D4 head-tail;
D5a structural cost; optional D6 churn; future D7 generator behavior

\vspace{2pt}\centering$\downarrow$\vspace{2pt}

\raggedright
\textbf{Outputs}: CSV/Markdown/\LaTeX{} audit tables and failure-case slices
\end{minipage}}
\caption{\tool artifact pipeline. The current resource implements D1--D5a over
item-to-\sid mappings, includes optional D6 churn support, and reserves D7 for
artifacts that expose generated candidates or beam traces.}
\Description{Pipeline diagram showing tokenizer exports normalized into
sid assignments, metadata, interactions, optional generator outputs, validation
checks, artifact diagnostics, and generated audit tables.}
\label{fig:audit-sid-pipeline}
\end{figure}

\paragraph{D1: utilization.}
\tool reports per-level usage, prefix counts, entropy-style imbalance
summaries, and dead or underused code indicators. These expose codebook
collapse and uneven capacity even when full ranking metrics are unavailable.

\paragraph{D2: collision profile.}
\tool measures duplicate full \sid{}s and prefix collisions at each depth. In
the current resource version, D2 is a collision profile rather than a causal
harm estimate. Interaction-qualified collision harm, as motivated by
qualification-aware collision work~\cite{hu2026quasid}, is reserved for the
next artifact version.

\paragraph{D3: semantic-collaborative alignment.}
Semantic grouping and collaborative usefulness need not agree. \tool therefore
computes a co-occurrence-based collaborative neighborhood proxy and asks how
often prefix neighborhoods recover those behaviorally related items. Metadata
purity is retained only as an auxiliary diagnostic, not as the definition of
collaborative quality.

\paragraph{D4: head-tail capacity.}
\tool splits items into popularity buckets and measures whether full \sid
capacity and prefix structure are allocated differently across head, mid, and
tail items. This helps separate genuine tokenizer behavior from simply
restating the catalog popularity distribution.

\paragraph{D5a, D6, and D7.}
D5a reports structural cost proxies: \sid length, prefix fan-out, duplicate
codes, and trie-like expansion pressure. These are not real generator serving
latencies. D6 computes churn between tokenizer refreshes and is useful for
drift-aware continual-tokenization settings~\cite{feng2026dact}. D7 is an
interface hook for generator behavior, such as invalid paths, next-token
entropy, and duplicate generated candidates; it requires
\texttt{generator\_outputs} and is not claimed by the current evidence.

Identifier foundations and recommender diagnostic frameworks motivate the
resource framing, but they are not runnable tokenizer rows. Table
\ref{tab:method-coverage} therefore focuses on SID artifact facets and marks
which parts are demonstrated in v0.

\begin{table*}[t]
\caption{Method/facet coverage in the current \tool resource. R0 foundation and
diagnostic-tooling citations motivate the problem but are not listed as
runnable tokenizer artifacts.}
\label{tab:method-coverage}
\scriptsize
\setlength{\tabcolsep}{3.5pt}
\begin{tabularx}{\textwidth}{@{}l l l l Y@{}}
\toprule
Facet & Anchor examples & v0 artifact role & Diag. & Boundary \\
\midrule
A canonical SID & GRID/RQ-KMeans & runnable main A & D1--D5a &
Musical row uses processed feature text, not raw-text TIGER. \\
B1 collaborative / predictability & ReSID/GAOQ & runnable main B & D1--D5a,D3 &
Bounded Musical export; not representative of all B facets. \\
B2 collision / capacity & QuaSID, AdaSID, CARD, CapsID & literature/backlog &
D1,D2,D4 & Motivates diagnostics; no faithful named-method result claimed. \\
B3 ranking / retrieval & DIGER, joint search-rec SID & interface target &
D3,D5a,D7 & D7 needs generated candidates or beam traces. \\
B4 bottleneck / interface & AsymRec, CapsID, Long SID & literature target &
D4,D5a,D7 & Current v0 has structural D5a only. \\
C drift / staleness & DACT, SID staleness & optional smoke & D6 &
Drift evidence, not a Cluster B replacement. \\
D deployment / search & Snapchat SID, search-rec SID & motivation target &
D5a,D7 & No production serving or online-impact claim. \\
Controls / portability & sanity, MovieLens smoke & runnable controls &
D1--D5a & Metric checks only; not SID methods. \\
\bottomrule
\end{tabularx}
\end{table*}
